\section{Learning Outcomes}

\subsection{Tasks Accomplished}
I completed manual implementations of both required models and their full training/evaluation pipelines in \texttt{MLP\_template.py} and \texttt{CNN\_template.py}. The main parts I finished were:
\begin{itemize}
	\item Manual cross-entropy loss with numerical-stability fixes (max-shift Softmax and an $\epsilon$ term in the logarithm).
	\item Manual backpropagation for both models. For MLP, I computed Sigmoid-layer derivatives and updated parameters from those gradients. For CNN, I propagated gradients from the dense output layer back through flattening and then into the convolution kernel.
	\item Debugging CNN tensor shapes so each step matched correctly. I checked that convolution outputs, flattened vectors, and reshaped gradients all stayed compatible during both forward and backward passes.
	\item Targeted hyperparameter checks across the values I benchmarked, including MLP hidden size, MLP initialization scale, and CNN kernel size, to justify the final settings used in the templates.
\end{itemize}

\subsection{Challenges and Debugging Realities}
The main challenge was debugging the manual training logic without autograd to catch mistakes for me. I had to check the forward pass, the gradients, and the parameter updates by hand, and the failed-guess training flowchart (Methodology, Figure \ref{fig:mlp_training_walkthrough}) was one of the investigative steps I used to confirm that the correction signal was behaving as expected.

The CNN side was also constrained by the fixed 5-epoch limit in the project. That meant there was little room for slow convergence, so I had to keep the implementation simple, check the tensor shapes carefully, and use kernel-size benchmarking to choose a setting that learned reasonably well within the short training budget.

\subsection{Lessons Learned}
These experiments helped me understand how architecture choices and training budget affect results:
\begin{itemize}
	\item The final MLP setting (hidden size 128, learning rate 0.5, 30 epochs) reached $0.9768$ test accuracy. Hidden-size testing showed negligible gains beyond 128 neurons (both 128 and 256 reached $0.967$), so that represents a practical capacity point for this setup.
	\item CNN test accuracy moved nonlinearly across epochs ($0.920 \rightarrow 0.915 \rightarrow 0.938 \rightarrow 0.947 \rightarrow 0.930$) even though training loss decreased consistently. This pattern reflects the very tight 5-epoch budget, which left little room for stable convergence.
	\item Kernel-size testing selected $7\times 7$ as the best option, achieving $0.933$ test accuracy under the 5-epoch constraint, confirming its selection in \texttt{CNN\_template.py}.
	\item CNN batch-size analysis showed batch 64 performed best ($0.9296$ final accuracy with $0.2673$ loss), validating the chosen batch size of 64 in the template.
	\item Under the required budgets, the final gap between MLP ($97.68\%$) and CNN ($92.96\%$) is $4.72$ percentage points. This gap reflects both the 6-fold longer training duration for MLP and the superior capacity of dense layers versus a single convolutional kernel for feature extraction on MNIST.
\end{itemize}
